\documentclass[11pt]{article}
\usepackage[margin=1in]{geometry}
\usepackage{amsmath,amssymb,amsthm,mathtools}
\usepackage{microtype}
\usepackage[hidelinks]{hyperref}

\newtheorem{theorem}{Theorem}
\newtheorem{corollary}[theorem]{Corollary}
\theoremstyle{remark}
\newtheorem*{remark}{Scope boundary}

\newcommand{\D}{\mathbb D}
\newcommand{\C}{\mathbb C}
\newcommand{\E}{\mathbb E}
\newcommand{\norm}[1]{\lVert #1\rVert}
\newcommand{\Zf}{\mathcal Z}

\title{An Iff Classification of Gaussian Power Series\
with the Peres--Vir\'ag Zero Process}
\author{Literature-implied resolution packet for arXiv:2103.11947}
\date{11 August 2026}

\begin{document}
\maketitle

\begin{center}
\fbox{\parbox{0.91\textwidth}{
\textbf{Classification: literature-implied answer; high confidence.}
Sodin's 2000 Calabi-rigidity theorem implies an exact converse to the
same-Bergman-process part of Mukeru--Mulaudzi.  A nondegenerate proper complex
Gaussian power series has the Peres--Vir\'ag zero process if and only if it is,
in law, a zero-free analytic multiple of the iid Gaussian series.  When its
coefficient covariance is a bounded invertible operator on $\ell^2$, this is
equivalent to the inverse covariance being Toeplitz.  The classification of
series whose zeros form an arbitrary determinantal process remains open.
}}
\end{center}

\section{Source problem and exact scope}

On PDF page 2 of \cite{MukeruMulaudzi}, the authors formulate the broad task of
determining dependent Gaussian coefficient sequences whose power-series zeros
are determinantal.  Their Theorem~1 on the same page proves a concrete
sufficient condition: if the coefficient covariance is $G^{-1}$ for a suitable
positive invertible Toeplitz matrix $G$, then the zeros have exactly the
Peres--Vir\'ag law, the determinantal process with Bergman kernel
\[
 L(z,w)=\frac{1}{\pi(1-z\overline w)^2}.
\]
The result below gives a complete converse for this same zero process.  Its
core is already contained in Sodin's Theorem~2 \cite{Sodin}, which appears on
PDF page 5 and predates the source paper by two decades.

Throughout, ``proper'' means zero pseudo-covariance.  Nondegeneracy means that
the analytic Gaussian law has no nontrivial deterministic linear relation; it
is the hypothesis corresponding to Sodin's linearly independent components.

\section{Exact classification}

Let
\[
 f(z)=\sum_{n\ge0}\xi_nz^n,
 \qquad
 g(z)=\sum_{n\ge0}\eta_nz^n,
\]
where $f$ is a nondegenerate centred proper complex Gaussian analytic function
on $\D$ and the $\eta_n$ are iid standard proper complex Gaussians.  Write
\[
 K_f(z,w)=\E[f(z)\overline{f(w)}],
 \qquad K_0(z,w)=\frac1{1-z\overline w}.
\]

\begin{theorem}[Peres--Vir\'ag iff classification]\label{thm:iff}
The following are equivalent.
\begin{enumerate}
\item The zero process of $f$ has the Peres--Vir\'ag law.
\item The expected zero measure of $f$ equals that of $g$, namely
\[
 \E n_f(dz)=\frac{1}{\pi(1-|z|^2)^2}\,dA(z).
\]
\item There is a zero-free analytic function $h$ on $\D$ such that
\[
 K_f(z,w)=\frac{h(z)\overline{h(w)}}{1-z\overline w}.
\]
\item As random analytic functions, $f\stackrel d=hg$.
\item If $h(z)=\sum_{j\ge0}h_jz^j$, then the coefficient sequence has the
same joint law as
\[
 \xi_n=\sum_{k=0}^{n}h_{n-k}\eta_k,\qquad n\ge0.
\]
\end{enumerate}
In particular, equality of the one-point intensity already forces equality of
the entire zero process.
\end{theorem}

\begin{proof}
The implication (1)$\Rightarrow$(2) is immediate.  By the
Edelman--Kostlan formula,
\[
 \E n_f=\frac{1}{4\pi}\Delta\log K_f(z,z)\,dA(z),
\]
with the equivalent $\partial\bar\partial$ normalization.  Thus (2) makes
\[
 \log K_f(z,z)-\log K_0(z,z)
\]
harmonic.  Since $\D$ is simply connected, it equals $\log|h(z)|^2$ for a
single-valued zero-free analytic $h$.  Hence
\[
 K_f(z,z)=|h(z)|^2K_0(z,z).
\]
The diagonal is a uniqueness set for kernels analytic in $z$ and
anti-analytic in $w$: expand at any point, or differentiate in $z$ and
$\overline w$.  Polarization therefore gives (3).  This is exactly the kernel
form of Sodin's Calabi-rigidity theorem.

A centred proper complex Gaussian analytic law is determined by its covariance
kernel.  The right-hand side in (3) is the covariance kernel of $hg$, proving
(4).  Comparing Taylor coefficients proves (5).  Conversely, (5) gives
$f\stackrel d=hg$; because $h$ has no zeros, multiplication by $h$ leaves every
zero unchanged.  Thus (5)$\Rightarrow$(1), completing the cycle.
\end{proof}

\section{The inverse-Toeplitz converse}

The preceding theorem needs no boundedness assumption on the infinite
coefficient covariance.  Under the standard bounded-operator hypothesis it
also characterizes the source's matrix condition.

\begin{corollary}[Toeplitz inverse iff]\label{cor:toeplitz}
Assume the coefficient covariance matrix $C=(\E\xi_j\overline{\xi_k})$ defines
a bounded positive invertible operator on $\ell^2(\mathbb N_0)$.  Then the five
conditions of Theorem~\ref{thm:iff} hold if and only if $C^{-1}$ is a bounded
positive invertible Toeplitz operator.
\end{corollary}

\begin{proof}
Under item (5), let $T_h$ be the lower-triangular analytic Toeplitz matrix of
multiplication by $h$ on $H^2$.  Then
\[
 C=T_hT_h^*.
\]
Boundedness and invertibility of $C$ force $T_h$ to be bounded and bijective;
equivalently $h,1/h\in H^\infty$.  Therefore
\[
 C^{-1}=T_{1/h}^*T_{1/h}=T_{|1/h|^2},
\]
which is Toeplitz.

Conversely, write $C^{-1}=T_\varphi$.  Positivity and bounded invertibility
give $0<c\le\varphi\le C_0$ almost everywhere.  Spectral factorization yields
an outer $a$ with $\varphi=|a|^2$ and $a,1/a\in H^\infty$.  Hence
\[
 C^{-1}=T_a^*T_a,
 \qquad
 C=T_{1/a}T_{1/a}^*.
\]
Taking $h=1/a$ gives item (3) of Theorem~\ref{thm:iff}.
\end{proof}

\begin{remark}
The source's introductory wording asks about any determinantal zero process.
The theorem above classifies exactly the process appearing in its main result:
the Bergman/Peres--Vir\'ag process.  For an unspecified determinantal kernel,
the one-point intensity does not encode the determinantal representation, and
the general Kac--Rice correlation identities lead to a substantially broader
functional classification problem.  No claim about that unrestricted problem
is made here.
\end{remark}

\paragraph{Literature status.}
Sodin's Theorem~2 says that two nondegenerate Gaussian analytic functions with
the same expected zero measure differ, at the analytic-vector level, by a
zero-free analytic scalar and a constant unitary.  Applying it to
$(1,z,z^2,\ldots)$ gives Theorem~\ref{thm:iff}.  The coefficient convolution
and Corollary~\ref{cor:toeplitz} are short operator-theoretic translations of
that established theorem.  The target paper does not cite Sodin.  Searches on
11 August 2026 found no paper explicitly recording this converse in the
target's coefficient-matrix language, so the result is classified as a
literature implication rather than a new proof claim.

\paragraph{Evidence locations.}
Source: arXiv:2103.11947, PDF page 2 (broad problem and Theorem~1), continuing
on page 3.  Supporting source: arXiv:math/0007030, Section~3, Theorem~2 on PDF
page 5; its proof continues on pages 5--6.

\begin{thebibliography}{9}
\bibitem{MukeruMulaudzi}
S.~Mukeru and M.~P. Mulaudzi,
\emph{Zeros of Gaussian Power Series, Hardy Spaces and Determinantal Point
Processes}, arXiv:2103.11947 (2021).

\bibitem{Sodin}
M.~Sodin,
\emph{Zeros of Gaussian Analytic Functions},
Math. Res. Lett. \textbf{7} (2000), 371--381;
arXiv:math/0007030.

\bibitem{PeresVirag}
Y.~Peres and B.~Vir\'ag,
\emph{Zeros of the I.I.D. Gaussian Power Series: A Conformally Invariant
Determinantal Process},
Acta Math. \textbf{194} (2005), 1--35.
\end{thebibliography}

\end{document}
